\paragraph{\trIsomorphism{}、\trDeformation。}\label{sub:isomorphisms-deformations}
两个\trFibrationWithSingularities{}是\trIsomorphic{}的，如果存在\trTotalSpace{}之间与底空间之间的（作为\trExactSymplectic\trCorneredBoundary\trManifold{}的）\trIsomorphism，它们与\trFibration{}映射\trCompatible。如何定义这类\trFibration{}的一个一般的\trDeformation{}概念，同样显而易见：即把所涉及的全部几何数据做成一个单参数\trFamily。我们并不这样做，而是把注意力集中在两种更特殊的\trDeformation{}上，它们将在后面若干技术性结果的表述中出现。第一种的主要特征是，\trRegular\trFiber{}的\trIsomorphism{}类得到保持；而在第二种中，\trFiber{}的结构完全不变。
取定\trManifold{}$E$ 与 $B$ 以及它们之间的一个映射 $\pi$。在其上取\trSmooth{}的\trExactSymplectic{}结构单参数\trFamily{}$(\phi^s_E,\theta^s_E,I^s_E)$、$(\phi^s_B,\theta^s_B,I^s_B)$，$s\in[0;1]$，使得对每个 $s$，$\pi$ 都成为一个\trFibrationWithSingularities。如果存在 $E$ 上\trSmooth{}的函数 $R^s$，使得由
\begin{equation}\label{eq:well-behaved-vector-field}
\phi^s_E(\,\cdot\,,Y^s)=\tfrac{\partial}{\partial s}\theta^s_E-dR^s
\end{equation}
所定义的\trVectorField{}$Y^s$ 都平行于 $\partial^h E$，我们就说这个\trDeformation{}沿 $\partial^h E$ 是\textbf{\trWellBehaved}的（well-behaved）。
一个\trEquivalence{}的表述是：取乘积 $\tilde{\pi}=\mathrm{Id}\times\pi\colon\widetilde{E}=[0;1]\times E\to\widetilde{B}=[0;1]\times B$，并在其\trTotalSpace{}上装备 $\phi_{\widetilde{E}}=d(\theta^s_E+R^s\,ds)$。这在 $\ker(D\tilde{\pi})$ 上显然是辛的，因此我们得到 $T\widetilde{E}$ 的一个如 \eqref{eq:tangent-splitting} 的\trSplitting；而 \eqref{eq:well-behaved-vector-field} 所说的是，该\trSplitting{}的\trHorizontal{}部分 $T\widetilde{E}^h$ 平行于 $\partial^h\widetilde{E}=[0;1]\times\partial^h E$。这使得人们可以为 $[0;1]\times(B\setminus\mathrm{Critv}(\pi))$ 中的路径定义\trParallelTransport{}映射。特别地，沿 $[0;s]\times\{y\}$ 行走，便得到\trExactSymplectic{}的\trIsomorphism
\begin{equation}\label{eq:deformation-fibre-isomorphism}
k^s_y\colon(E_y,\phi^0_{E_y},\theta^0_{E_y})\longrightarrow(E_y,\phi^s_{E_y},\theta^s_{E_y}).
\end{equation}
更具体地说，这些\trIsomorphism{}是这样定义的：积分 $Y^s|_{E_y}$ 的（相对于 $\phi^s_E$ 的）\trVertical{}分量——如果人们取 $\partial_s\theta^s_{E_y}-dR^s|_{E_y}$ 作为 $\partial_s\phi^s_{E_y}$ 的原函数，这将是\trMoser{}论证中的那个\trVectorField。诸 $k^s_y$ \trSmooth{}地依赖于 $(s,y)$，并且在相差\trHamiltonianIsotopy{}的意义下，与沿 $B\setminus\mathrm{Critv}(\pi)$ 中路径的\trParallelTransport{}交换。换言之，$E$ 的\trRegular{}部分的\trFiberBundle{}结构，在沿 $\partial^h E$ \trWellBehaved{}的\trDeformation{}之下保持不变。

另一个限制性更强的概念，是\textbf{\trConstantAlongFibres}（constant along the fibres）的\trDeformation。它由与前面同样的单参数\trFamily{}$(\phi^s_E,\theta^s_E,I^s_E)$ 与 $(\phi^s_B,\theta^s_B,I^s_B)$ 给出，只是现在我们还要求 1-\trForm{}$\partial_s\theta^s_E-dR^s$ 在 $TE^v$ 上为零。显然，这意味着 \eqref{eq:deformation-fibre-isomorphism} 的那些映射都是恒同。除此之外，我们还要求在 $\partial^h E$ 的一个邻域内，$I^s_E|_{TE^v}$ 与 $s$ 无关。

设 $\pi\colon E\to S$ 是一个\trFibrationWithSingularities，其中 $S$ 是一个带边界的\trSurface。设 $\kappa$ 是 $E$ 上一个在 $TE^v$-方向上消失的 1-\trForm，并且在 $\partial E\cup\mathrm{Crit}(\pi)$ 的一个邻域内为零。对 $s\in[0;1]$ 与某个常数 $c\gg 0$，考虑
\begin{equation}\label{eq:kappa-deformation-family}
\begin{aligned}
&\theta^s_S=(cs+1)\theta_S,\qquad \phi^s_S=d\theta^s_S,\qquad I^s_S=I_S,\\
&\theta^s_E=\theta_E+s\cdot\kappa+cs\cdot\pi^*\phi_S,\qquad \phi^s_E=d\theta^s_E,
\end{aligned}
\end{equation}

由于 $\omega^s_E|_{TE^v}=\omega_E|_{TE^v}$ 是\trNondegenerate{}的，我们得到一个依赖于 $s$ 的\trSplitting{}$TE_x=TE^{h,s}_x\oplus TE^v_x$，它是 \eqref{eq:tangent-splitting} 的\trDeformation。通过选取足够大的 $c$，我们可以保证 $\phi^s_E|_{TE^{h,s}_x}$ 是正的（相对于从底空间诱导的\trOrientation），从而 $\phi^s_E$ 是辛的。取 $I^s_E$ 为唯一的与 $\phi^s_E$ \trCompatible{}的\trAlmostComplex{}结构，它在 $TE^v$-方向上等于 $I_E$，而在 $TE^{h,s}$ 上与 $I_S$ 的\trPullback{}一致。由于在 $\partial E$ 附近 $TE^{h,s}=TE^h$，$I^s_E$ 在那里将与 $I_E$ 一致，这就蕴含边界是弱 $I^s_E$-\trConvex{}的。其次，与 $\theta^s_E$ 和 $\phi^s_E$ 相伴的负\trLiouville\trVectorField{}沿 $\partial^h E$ 严格\trIncoming，因为这条性质只依赖于 $TE^v$ 加项；参见\cref{thm:fibre-exact-symplectic}的证明。现在取一条\trBoundaryCircle{}$C\subset\partial S$，以及相应的 $F=\pi^{-1}(C)\subset\partial^v E$。对 $x\in F$，$TE^h_x$ 包含一个一维子空间 $TF_x\cap TE^h_x$，而且通过\trProjection{}到 $S$ 并使用那里的边界\trOrientation，它还有一个\trCanonical{}的\trOrientation。负\trLiouville\trVectorField{}沿 $F$ \trIncoming{}这一事实，\trEquivalence{}于
\begin{equation}\label{eq:liouville-inwards-criterion}
\theta_E\,|\,(TF_x\cap TE^h_x)>0.
\end{equation}

由于 $\theta_S$ 本身使 $S$ 成为一个\trExactSymplectic\trCorneredBoundary\trManifold，故有 $\pi^*\theta_S\,|\,(TF_x\cap TE^h_x)>0$。这一点，加上在 $\partial^v E$ 附近 $\theta^s_E=\theta_E+cs\cdot\pi^*\phi_S$ 的事实，蕴含 \eqref{eq:liouville-inwards-criterion} 的类似物在整个\trDeformation{}过程中继续成立。至此我们已经验证了 $(E,\phi^s_E,\theta^s_E,I^s_E)$ 对每个 $s$ 都是\trExactSymplectic\trCorneredBoundary\trManifold，从而 \eqref{eq:kappa-deformation-family} 是一个\trConstantAlongFibres{}的\trDeformation。从\trFiberBundle{}的观点看，其效果是把\trSymplectic\trConnection{}改变 $\kappa$。为将来使用，我们记录一个具体的应用：
\begin{lemma}\label{thm:transport-isotopy-realization}
设 $\pi\colon E\to S$ 是一个\trFibrationWithSingularities，其中 $S$ 是一个带边界的\trSurface。设 $\beta\colon[0;1]\to\mathrm{int}(S)\setminus\mathrm{Critv}(\pi)$ 是一条\trEmbedding{}路径，而 $(\phi^s)$ 是 $\mathrm{Aut}(E_{\beta(0)},\partial E_{\beta(0)})$\footnote{译注：原书此处印作 $E_{c(0)}$、$\partial E_{c(0)}$；疑应为 $E_{\beta(0)}$、$\partial E_{\beta(0)}$（$c$ 在前文仅作常数 $c\gg 0$ 出现，与此处语义不合，且惟 $\beta(0)$ 使 $h^s_\beta=h_\beta\circ\phi^s$ 有意义）；此处径改。——译者注} 中以 $\phi^0=\mathrm{id}$ 开始的一个\trIsotopy。那么存在一个\trConstantAlongFibres{}的\trDeformation，它在 $s=0$ 处从给定的\trFibration{}结构出发，使得沿 $\beta$ 所得的\trParallelTransport{}映射族满足
\begin{equation*}
h^s_\beta=h_\beta\circ\phi^s.
\end{equation*}
此外，若 $U\subset S$ 是一个包含 $\beta$ 的一部分的\trOpen{}子集，则可以把这个\trDeformation{}安排成使 $\partial_s(\theta^s_E-\pi^*\theta^s_S)$ 在 $\pi^{-1}(U)$ 之外为零。
\noqed
\end{lemma}
